\section{Analiza wydajności Monte Carlo}

\subsection{Konfiguracje testowe}

Przeprowadzono eksperymenty dla czterech konfiguracji równoległych (łącznie 8 rdzeni):
\begin{itemize}
    \item \textbf{n1c8}: 1 węzeł × 8 wątków OpenMP (czysty OpenMP)
    \item \textbf{n2c4}: 2 węzły × 4 wątki (hybrid)
    \item \textbf{n4c2}: 4 węzły × 2 wątki (hybrid)
    \item \textbf{n8c1}: 8 węzłów × 1 wątek (czysty MPI)
\end{itemize}

Testowano symulację Monte Carlo do szacowania wartości $\pi$ dla liczby punktów od 1 miliarda do 30 miliardów.

\subsection{Obserwacje}

Głównym odkryciem eksperymentów jest \textbf{brak znaczących różnic w wydajności} między konfiguracjami:

\begin{itemize}
    \item \textbf{n1c8} (1 węzeł, 8 wątków OpenMP): 624.2--625.8 Mpoints/s
    \item \textbf{n2c4} (2 węzły, 4 wątki): 615.7--624.8 Mpoints/s
    \item \textbf{n4c2} (4 węzły, 2 wątki): 613.7--624.6 Mpoints/s
    \item \textbf{n8c1} (8 węzłów, czysty MPI): 605.1--624.3 Mpoints/s
\end{itemize}

\textbf{Różnice między konfiguracjami wynoszą < 1.5\%}, co jest w granicach szumu pomiarowego i wariancji systemu.

\paragraph{Stabilność wydajności.} Niezależnie od rozmiaru problemu (1-30 miliardów punktów), wszystkie konfiguracje utrzymują praktycznie stałą przepustowość \textasciitilde{}623 Mpoints/s. Nie występuje degradacja wydajności wraz ze wzrostem rozmiaru danych.

\paragraph{Niewrażliwość na topologię.} Czy obliczenia odbywają się:
\begin{itemize}
    \item w jednym węźle (n1c8)
    \item w dwóch węzłach połączonych siecią (n2c4)
    \item w ośmiu węzłach z komunikacją międzywęzłową (n8c1)
\end{itemize}
nie ma praktycznie żadnego znaczenia dla wydajności końcowej.

\paragraph{Idealna skalowalność.} Dla 8 rdzeni procesora uzyskujemy \textasciitilde{}625 Mpoints/s, co oznacza \textasciitilde{}78 Mpoints/s na rdzeń -- niezależnie od konfiguracji. Jest to niemal idealny speedup $S_8 \approx 8$.

\subsection{Wyjaśnienie zjawiska}

Monte Carlo to \textbf{embarrassingly parallel problem} -- obliczenia są całkowicie niezależne:

\begin{lstlisting}[language=C]
#pragma omp parallel
{
    unsigned int seed = (unsigned int)time(NULL) ^ omp_get_thread_num();
    long local_count = 0;

    #pragma omp for
    for (long i = 0; i < num_points; i++) {
        double x = (double)rand_r(&seed) / RAND_MAX;  // Niezależny seed
        double y = (double)rand_r(&seed) / RAND_MAX;
        if (x*x + y*y <= 1.0) local_count++;
    }

    #pragma omp atomic
    count += local_count;
}
\end{lstlisting}

\subsubsection{Brak cache coherence overhead}

W przeciwieństwie do N-body, gdzie wszystkie wątki czytały współdzielone tablice pozycji, w Monte Carlo:
\begin{itemize}
    \item Każdy wątek/proces generuje \textbf{własne losowe punkty}
    \item Brak współdzielonych danych w pętli głównej
    \item Zero ruchu w protokole spójności cache
\end{itemize}

\subsubsection{Minimalny overhead komunikacji}

Komunikacja występuje tylko \textbf{raz na końcu} obliczeń:

\begin{lstlisting}[language=C]
// Każdy proces oblicza lokalnie wartość Pi
double local_pi = monte_carlo_pi(points_per_rank);

// MPI: redukcja jednorazowo po zakończeniu wszystkich obliczeń
MPI_Reduce(&local_pi, &global_pi, 1, MPI_DOUBLE, 
           MPI_SUM, 0, MPI_COMM_WORLD);
global_pi /= size;  // średnia
\end{lstlisting}

Dla 30 miliardów punktów i czasu wykonania \textasciitilde48s, koszt pojedynczej operacji \texttt{MPI\_Reduce} (kilka $\mu$s) jest zaniedbywalny.

\subsubsection{Idealna skalowalność teoretyczna}

Dla problemu embarrassingly parallel, przyspieszenie teoretyczne wynosi:
\begin{equation}
    S_p = p
\end{equation}
gdzie $p$ to liczba procesorów. Wszystkie konfiguracje wykorzystują 8 rdzeni, więc teoretycznie powinny osiągnąć tę samą wydajność -- co dokładnie obserwujemy.

\subsubsection{Dlaczego n1c8 jest minimalnie szybsze}

Czysty OpenMP (n1c8) ma przewagę ~1\% z powodu:
\begin{itemize}
    \item \textbf{Brak overhead'u startowego MPI}: inicjalizacja \texttt{MPI\_Init} i finalizacja
    \item \textbf{Szybsza redukcja}: OpenMP \texttt{reduction(+:count)} w pamięci współdzielonej vs \texttt{MPI\_Reduce} przez sieć
    \item \textbf{Brak synchronizacji międzywęzłowej}: wszystko dzieje się w jednym węźle
\end{itemize}

Jednak w praktyce różnica jest na poziomie szumu pomiarowego.

\subsection{Wnioski}

\begin{enumerate}
    \item Dla problemów \textbf{embarrassingly parallel} wybór architektury (OpenMP vs MPI) ma \textbf{znikomy wpływ} na wydajność (~1\% różnicy).
    
    \item Wszystkie konfiguracje osiągają praktycznie \textbf{idealną skalowalność}:
    \begin{equation}
        S_p \approx p
    \end{equation}
    gdzie $p=8$ (liczba rdzeni). Wydajność jest proporcjonalna do liczby rdzeni, niezależnie od topologii.
    
    \item Kluczowe czynniki zapewniające dobrą skalowalność:
    \begin{itemize}
        \item \textbf{Brak współdzielonych danych} -- każdy wątek/proces operuje niezależnie
        \item \textbf{Minimalna komunikacja} -- jedna operacja redukcji stanowi $< 0.001\%$ czasu wykonania
        \item \textbf{Równomierne obciążenie} -- każdy rdzeń przetwarza tę samą liczbę punktów
    \end{itemize}
    
    \item Nieznaczna przewaga czystego OpenMP (n1c8) wynika z:
    \begin{itemize}
        \item Braku overhead'u inicjalizacji/finalizacji MPI
        \item Szybszej redukcji w pamięci współdzielonej (\texttt{\#pragma omp atomic})
        \item Braku synchronizacji międzywęzłowej
    \end{itemize}
    
    \item \textbf{Rekomendacja praktyczna}:
    \begin{itemize}
        \item Dla pojedynczego węzła: użyj \textbf{czystego OpenMP} (n1c8) -- prostsze, równie wydajne
        \item Dla wielu węzłów: użyj \textbf{czystego MPI} (n8c1) -- gdy problem wymaga więcej pamięci niż jeden węzeł
        \item Konfiguracje hybrydowe nie dają dodatkowych korzyści dla tego typu problemu
    \end{itemize}
    
    \item Problem stanowi \textbf{wzorcowy przykład idealnej równoległości} -- brak interakcji między obliczeniami, minimalna synchronizacja, liniowa skalowalność.
\end{enumerate}
